\section{} % 2

\subsection{} % 2.1
The following is the full algorithm implementation of the Procrustes algorithm:

\begin{lstlisting}
def procrustes(S: NDArray[Any], target: NDArray[Any]) -> NDArray[Any]:
    M, d, N = S.shape
    assert d == 2, "Only 2D coordinates are supported in this implementation."
    assert target.shape == (2, N), (
        f"Target shape must be (2, N)={(2, N)}, but got {target.shape}."
    )

    ### Translate all coordinates to origo
    x_center = S[:, 0, :].mean(axis=1, keepdims=True)  # shape (182, 1)
    y_center = S[:, 1, :].mean(axis=1, keepdims=True)  # shape (182, 1)

    centered = S.copy()
    centered[:, 0, :] -= x_center  # broadcasts across 32 landmarks
    centered[:, 1, :] -= y_center
    target -= target[0].mean(axis=0)  # Center the target as well
    target -= target[1].mean(axis=0)  # Center the target as well

    ### Scale all of them w.r.t. target
    enum = np.sum(centered * target, axis=(1, 2))
    denom = np.sum(centered * centered, axis=(1, 2))
    # Check for 0 division
    if np.any(denom == 0):
        print(
            "WARNING: Some shapes have zero variance. "
            "To avoid division by zero, 0-s are set to 1.0"
        )
    denom = np.where(denom == 0, 1.0, denom)  # avoid division by zero
    s = enum / denom  # shape (M,)

    ### Rotate them all w.r.t. target
    output = centered * s[:, np.newaxis, np.newaxis]  # broadcast over (2, 32)
    for i in range(M):
        other = output[i]
        R = get_rotation(target, other)
        output[i] = R @ other

    return output
\end{lstlisting}
where the \lstinline{get_rotation()} function is as following:
% ROTATION FUNCTION
\begin{lstlisting}
def get_rotation(target: NDArray[Any], other: NDArray[Any]) -> NDArray[Any]:
    A = other @ target.T
    U, _, Vt = np.linalg.svd(A)
    V = Vt.T
    return V @ U.T
\end{lstlisting}

This particular implementation is limited to aligning $2D$ shapes.
The function asserts correct inputs, applies the centering by leveraging \texttt{numpy} broadcasting,
and similarly scales by leveraging broadcasting. 
Essentially the function computes the sum of the dot products as prescribed by the algorithm,
but does so by element-wise multiplication and summation across the appropriate axes.
The result is an $M$-dimensional vector with $M$ scaling factors, one for each shape.
Finally we apply rotation on every input, by computing the optimal rotation matrix using the 
SVD of the cross-covariance matrix between the target and the scaled shape.

The resulting aligned shapes are shown in \autoref{fig:10wings}:

\begin{figure}[H]
	\centering
	\includegraphics[width=0.8\textwidth]{ass8/10_wings.png}
	\caption{10 Wing Shapes given by 32 landmarks, aligned using the Procrustes algorithm.}
	\label{fig:10wings}
\end{figure}
which demonstrates that the algorithm successfully aligns the shapes.



\subsection{} % 2.2

The (favorite) classifier chosen is a Support Vector Machine (SVM).
Technically the SVM provides optimal separating hyperplane between classes of data points,
and by evaluating the sign of the decision function for a given point relative to that boundary
it becomes a classifier.

For this particular task, with $2\times 32$ landmarks making up $n=64$ features,
the choice of kernel seems to be a linear kernel.
Typically, with geometric data such as this, which have such concrete geometric meaning,
we expect distinct races to have inter-race linearly seperable distinctions.

We achieve the following results:
\begin{itemize}
	\item Unaligned Accuracy: $0.8228$
	\item Aligned Accuracy: $0.9367$
\end{itemize}
which illustrates a clear improvement in performance.
This is likely due to Procrustes alignment serving as a form of data normalization 
i.e. alignment in terms of center, scale, and orientation.

This makes the classifier more robust to e.g. outliers, which may have completely different orientation,
or overblown scale. What we saw in figure \ref{fig:10wings} wasn't directly an outlier,
but indeed we saw the inpact of alignment. The accuracy result seems to reflect, that such alignment
makes it easier for the SVM with linear kernel to find 15 separating hyperplanes to distinguish between the 15 different species.
\footnote{Each hyperplane is a One-vs-Rest hyperplane, making up 15 in total.}

\subsection{} % 2.3

We very concretely removed the spatial orientation, size, and position information when performing Procrustes alignment on the dataset.
What we kept are the relative spatial relationships between the 32 landmarks.
The assumption is, that what matters is the shape. Its particular features and characteristics should be rotation, scale, and translation invariant.

Even though we saw improved accuracy performance between unaligned and aligned training data onto test data,
without better domain knowledge we can't be sure that e.g. size plays a crucial role in determining the species.
The improved performance in terms of accuracy led us to believe that these alignments removed
some factors that acted as noise rather than signal, in order to classify the species.
However it is important to be cognisant of what we are giving up, and consider whether these
could in fact be relevant or even crucial in solving the task at hand e.g. species classification.
